\chapter{Validation} \label{chap: Validation}

\section{Smooth wall pipe}\label{sec: Smooth wall pipe}


\section{Sinusoidal wall pipe}\label{sec: Sinusoidal wall pipe}

\subsection{Postprocessing rough Simulations}\label{subsec: postprocess rough Simulations}

In order to obtain correct data from the rough simulations, the friction calculation in the postprocessing must be expanded.
The friction calculation is changed, such it is able to account for the rough surface.\\
DESCRIBE NEW POSTPROCESSING 
\\
To be able to validate the friction results of the \Name{Python}-postprocessing script, $\tau_w$ is computed by using the force scaling algorithm of \Name{Neko}.
The force scaling factor $alpha$, that is used in \Name{Neko} to ensure a constant mass flux, can also be used to compute $Re_\tau$ with
\begin{equation}
    Re_\tau = \sqrt{\alpha*0.5}*Re_b
\end{equation}
while $\alpha$ is computed using
\begin{equation}
    \alpha = \frac{C_{\text{target}} - C_{\text{current}}}{C_{\text{base}}}
\end{equation}
with $C_{\text{target}}$ denoting the desired mass flux, $C_{\text{current}}$ denoting the mass flux at the current time step and $C_{\text{base}}$ denoting the base mass flux, at the start of the simulation CHECKEN OB DAS mit C base STIMMT!!!.
